% Generated from locale/tr/content/applied-modal-logic/temporal-logic/introduction.tex; do not hand-edit.


\olfileid[tr]{aml}{tl}{int}

\olsection{Giriş}

Zamansal mantıklar, gerçekleşecek ya da gerçekleşmiş şeyler hakkındaki
iddiaları ele alır. Zamansal mantığın kurucusu olarak kabul edilen Arthur
Prior, buna \emph{zaman kipleri mantığı} adını vermiştir. Buradaki
incelememiz, büyük ölçüde Prior'ın zamansal işleçleri önerme mantığının temel
çerçevesine kattığı özgün modal yaklaşımı izleyecektir. Önerme mantığı ise
iddiaları genel olarak zaman kipinden yoksun sayar.

Örneğin önerme mantığında, köpeklerin âdeti olduğu üzere bazen oturan bazen
de oturmayan Beezie adlı bir köpekten söz edebilirim. Beezie'nin hem
oturduğunu hem de oturmadığını ileri sürmek klasik mantıkta çelişkili olurdu.
Fakat bu iki iddia elbette doğru olabilir; yalnızca aynı anda doğru olamazlar.
Dile zamansal işleçler eklemek bu iddiayı görece kolay biçimde dile
getirmemizi sağlar. Zamansal işleçlerin eklenmesi, ``Beezie bir ödül maması
ya da bir top alacak'' öncülünden ``Beezie bir ödül maması alacak ya da
Beezie bir top alacak'' sonucuna giden çıkarımların geçerliliğini de
açıklamamıza olanak verir.

Bununla birlikte zamansal mantık, bizi bir zamansal mantık çerçevesini
diğerine yeğlemeye götürebilecek birçok felsefi sorun doğurur. Örneğin
geleceğe ilişkin olumsal bir önerme, gelecek hakkında ne zorunlu ne de
imkânsız olan bir ifadedir. ``Richard yarın markete gidecek'' dediğimizde
henüz gerçekleşmemiş ve doğruluk değeri tartışmalı bir şey hakkında iddiada
bulunuruz. Aslında bu iddiaya daha baştan bir doğruluk değeri
\emph{atanıp atanamayacağı} bile tartışmalıdır. Katı belirlenimcilersek, söz
konusu olayın gerçekleşmesi beklenen zamandan önce bile bu tümcenin gerçekten
doğru ya da yanlış olduğu düşüncesini benimseyebiliriz; yalnızca doğruluk
değerini henüz bilmiyor olabiliriz. Buna karşılık geleceğe ilişkin olumsal
önermelerin doğruluk değerlerinin belirlenmemiş olduğu gerçekten açık bir
geleceğe inanabiliriz.

Zamanın yapısı ve doğası hakkındaki bu kabullerin birçoğunun, zamansal mantık
modelleri ve çerçeveleri arasındaki seçimlerimize yerleşmiş olduğu ortaya
çıkar. Örneğin zamanın doğrusal, dallanan, hatta döngüsel olduğu modeller
kurup kurmamamız gerektiğini sorabiliriz. Zamansal modellerimizin başlangıç
ve bitiş noktaları bulunup bulunmayacağına ve zamanın ayrık anlarla mı yoksa
bir süreklilik olarak mı temsil edileceğine karar vermemiz gerekebilir.

